\chapter{Advanced Hypothesis Testing, LRT \& Goodness-of-Fit}

In practical data science applications, population variance $\sigma^2$ is almost never known, requiring \textbf{Student's $t$-test}. Additionally, we frequently compare two populations or test if categorical data fits a theoretical model. In this final capstone chapter, we master \textbf{One-Sample and Two-Sample $t$-Tests}, \textbf{Chi-Squared ($\chi^2$) Goodness-of-Fit Tests}, and the general \textbf{Likelihood Ratio Test (LRT)}.

\section{Student's $t$-Test for Variance Unknown}

When $\sigma^2$ is unknown, we substitute the sample variance $S^2 = \frac{1}{n-1} \sum_{i=1}^n (X_i - \bar{X})^2$.

\begin{definitionbox}{One-Sample $t$-Test Statistic}
For $H_0: \mu = \mu_0$ with unknown $\sigma^2$:
$$T_{\text{calc}} = \frac{\bar{X}_n - \mu_0}{S / \sqrt{n}} \sim t_{n-1} \quad (\text{Student's } t \text{-distribution with } n-1 \text{ degrees of freedom})$$
The $t$-distribution has heavier tails than the standard normal $\mathcal{N}(0,1)$, accounting for additional uncertainty in estimating $S$.
\end{definitionbox}

\begin{videocard}{IITM BS Lecture: t-Test for Mean (Variance Unknown)}{https://www.youtube.com/watch?v=dbJwXIoDBtQ}
Official lecture deriving Student's $t$-test statistic, degrees of freedom, and critical value tables.
\end{videocard}

\section{Two-Sample Tests}

To test if two population means are equal ($H_0: \mu_1 - \mu_2 = 0$):

\begin{formulabox}{Two-Sample Pooled $t$-Test Formula}
Assuming equal variances $\sigma_1^2 = \sigma_2^2$:
$$T_{\text{calc}} = \frac{\bar{X}_1 - \bar{X}_2}{S_p \sqrt{\frac{1}{n_1} + \frac{1}{n_2}}} \sim t_{n_1 + n_2 - 2}$$
where pooled variance $S_p^2 = \frac{(n_1 - 1)S_1^2 + (n_2 - 1)S_2^2}{n_1 + n_2 - 2}$.
\end{formulabox}

\begin{videocard}{IITM BS Lecture: Two Samples from Normal Distribution}{https://www.youtube.com/watch?v=G4mJrRem0H0}
Official lecture deriving two-sample $z$-tests, pooled $t$-tests, and $F$-tests for variance equality.
\end{videocard}

\section{Likelihood Ratio Test (LRT) and Chi-Squared Goodness of Fit}

\begin{definitionbox}{Likelihood Ratio Test (LRT)}
For testing $H_0: \theta \in \Theta_0$ vs $H_1: \theta \in \Theta \setminus \Theta_0$:
$$\Lambda(\mathbf{x}) = \frac{\sup_{\theta \in \Theta_0} L(\theta)}{\sup_{\theta \in \Theta} L(\theta)}$$
By Wilks' Theorem, $-2 \ln \Lambda(\mathbf{x}) \xrightarrow{d} \chi_k^2$ where $k = \dim(\Theta) - \dim(\Theta_0)$.
\end{definitionbox}

\begin{formulabox}{Pearson's Chi-Squared ($\chi^2$) Goodness-of-Fit Test}
To test whether categorical frequency data $O_1, \dots, O_k$ matches expected counts $E_1, \dots, E_k$:
$$\chi_{\text{calc}}^2 = \sum_{i=1}^k \frac{(O_i - E_i)^2}{E_i} \sim \chi_{k - 1 - m}^2$$
Reject $H_0$ if $\chi_{\text{calc}}^2 \ge \chi_{\alpha, k-1-m}^2$.
\end{formulabox}

\index{t-Test}\index{Chi-Squared Test}\index{Likelihood Ratio Test}

\begin{lstlisting}[caption={SciPy/Python Implementation of t-Tests and Chi-Squared Goodness of Fit}]
import scipy.stats as stats
import numpy as np

# 1. One-Sample t-Test (H0: mu = 100)
data = np.array([102, 105, 98, 108, 101, 104, 99, 106, 103, 100, 107, 102, 105, 101, 104, 98])
t_stat, p_val = stats.ttest_1samp(data, popmean=100)
print(f"One-Sample t-test: t = {t_stat:.3f}, p-value = {p_val:.4f}")

# 2. Two-Sample Independent t-Test (Pooled)
group1 = [14, 16, 15, 13, 17, 15, 14, 16, 15, 15]
group2 = [11, 13, 12, 10, 14, 12, 13, 11, 12, 12]
t_ind, p_ind = stats.ttest_ind(group1, group2, equal_var=True)
print(f"Two-Sample t-test: t = {t_ind:.3f}, p-value = {p_ind:.4f}")

# 3. Chi-Squared Goodness-of-Fit Test (Die Fairness)
observed = np.array([15, 7, 8, 12, 10, 8])
expected = np.array([10, 10, 10, 10, 10, 10])
chi2_stat, p_chi2 = stats.chisquare(f_obs=observed, f_exp=expected)
print(f"Chi-Squared test: chi2 = {chi2_stat:.3f}, p-value = {p_chi2:.4f}")
\end{lstlisting}

\begin{videocard}{IITM BS Lecture: Likelihood Ratio Tests}{https://www.youtube.com/watch?v=KzXltUrRt84}
Official lecture deriving likelihood ratios and Wilks' asymptotic $\chi^2$ theorem.
\end{videocard}

\begin{videocard}{IITM BS Lecture: Goodness of Fit for Discrete Distributions}{https://www.youtube.com/watch?v=bnBk6mMEArU}
Official lecture computing observed vs expected cell counts and applying $\chi^2$ test statistic.
\end{videocard}


\newpage
\section{Practice Exercises}
\textit{Detailed step-by-step solutions are available in the Appendix.}

\begin{enumerate}
\item \textbf{One-Sample $t$-Test Calculation:} A sample of $n = 16$ measurements has sample mean $\bar{X} = 104$ and sample standard deviation $S = 8$. We test $H_0: \mu = 100$ vs $H_1: \mu \neq 100$.
    \begin{enumerate}
        \item State degrees of freedom $df$.
        \item Compute the test statistic $T_{\text{calc}}$.
        \item At $\alpha = 0.05$ with critical value $t_{0.025, 15} = 2.131$, state your decision.
    \end{enumerate}

\item \textbf{Two-Sample Pooled $t$-Test:} Two independent samples yield:
\begin{itemize}
    \item Group 1: $n_1 = 10, \bar{X}_1 = 15, S_1^2 = 4$.
    \item Group 2: $n_2 = 10, \bar{X}_2 = 12, S_2^2 = 6$.
\end{itemize}
    Compute pooled variance $S_p^2$ and $T_{\text{calc}}$ for $H_0: \mu_1 - \mu_2 = 0$.

\item \textbf{Chi-Squared Goodness of Fit (Die Fairness Test):} A 6-sided die is rolled $N = 60$ times, observing counts:
    $$\text{Observed } O_i = [15, 7, 8, 12, 10, 8] \quad \text{for faces } 1 \dots 6$$
    Under $H_0$: Die is fair ($E_i = 10$ for all $i$). Compute $\chi_{\text{calc}}^2$. At $\alpha = 0.05$ with critical value $\chi_{0.05, 5}^2 = 11.07$, is the die fair?

\item \textbf{Wilks' Likelihood Ratio Test Stat:} A likelihood ratio test yields likelihood ratio $\Lambda = 0.05$. Compute Wilks' test statistic $W = -2 \ln \Lambda$.

\item \textbf{Comprehensive Course Review Question ($\chi^2$ Independence Test):} In a marketing survey, 200 users are classified by Subscription Tier (Free, Premium) and Churn Status (Retained, Churned):
\begin{center}
\begin{tabular}{c|cc|c}
& Retained & Churned & Total \\
\hline
Free Tier & 80 ($E=90$) & 40 ($E=30$) & 120 \\
Premium Tier & 70 ($E=60$) & 10 ($E=20$) & 80 \\
\hline
Total & 150 & 50 & 200 \\
\end{tabular}
\end{center}
Compute the $\chi_{\text{calc}}^2$ test statistic using the given expected counts $E_{ij}$ to test for independence between Subscription Tier and Churn.
\end{enumerate}